\VerdierPhysicalPageAnchor{0036}%
\VerdierSemanticAnchor{ch1-1-1-2-degree-e-morphism-cont}
of degree $e$ satisfying the usual compatibility properties.

\VerdierSemanticAnchor{ch1-1-1-3}
\noindent\textbf{Definition 1.1.3.}\quad
A $G$-graded category is called \emph{additive} if its fiber category
$\mathcal C_e$ is additive [20].
\VerdierSemanticAnchor{ch1-1-1-3-additive-functor}
A graded functor between additive graded categories is called
\emph{additive} if its restriction to the fibers is additive.

\VerdierSemanticAnchor{ch1-1-1-3-cadd}
Let $\mathcal C$ be an additive $G$-graded category. The category
$\mathcal C$ with its fibered structure forgotten \emph{is not additive}
when $G\neq\{e\}$. One can nevertheless associate an additive category to
$\mathcal C$ as follows: for every pair $X,Y$ of objects of $\mathcal C$
and every element $g$ of $G$, the set
$\operatorname{Hom}_{\mathcal C}^{g}(X,Y)$ is naturally endowed with the
structure of an abelian group. Define a new category
$\mathcal C_{\mathrm{add}}$ by taking the objects of $\mathcal C$ as its
set of objects, and by taking as the set of morphisms between two objects
$X$ and $Y$:

\VerdierSemanticAnchor{ch1-1-1-3-eq1}
\begingroup
\makeatletter
\tagsleft@true
\let\veqno\@@leqno
\makeatother
\begin{equation}
  \operatorname{Hom}_{\mathcal C_{\mathrm{add}}}(X,Y)
  =\bigoplus_{g\in G}\operatorname{Hom}_{\mathcal C}^{g}(X,Y).
  \tag{1.1.3.1}
\end{equation}
\endgroup

\VerdierSemanticAnchor{ch1-1-1-3-cadd-composition}
Composition of morphisms in $\mathcal C_{\mathrm{add}}$ is defined in the
evident way, and one checks that $\mathcal C_{\mathrm{add}}$ is an
additive category.
\VerdierSemanticAnchor{ch1-1-1-3-fadd}
Likewise, given an additive cartesian functor
$F\colon\mathcal C\to\mathcal C'$ between two additive graded categories,
the same method constructs an additive functor
$F_{\mathrm{add}}\colon\mathcal C_{\mathrm{add}}\to
\mathcal C'_{\mathrm{add}}$. We shall never use this construction.

\VerdierSemanticAnchor{ch1-1-1-3-restrictive}
Observe that the definition of graded categories adopted here is more
restrictive than the one that naturally comes to mind. In addition to the
evident structures (1.1.2.1) and (1.1.2.2), we impose property (1.1.2.3).

\VerdierSemanticAnchor{ch1-1-2}
\subsection{Action of a Group on a Category up to Isomorphism}

\VerdierSemanticAnchor{ch1-1-2-1}
\noindent\textbf{1.2.1.}\quad
Let $\mathcal C$ be a $G$-graded category. Choose a normalized cleavage of
the fibered category $\mathcal C\xrightarrow{\deg}\mathcal G$ [18],
\emph{i.e.} for every $g\in G$ choose a base-change functor:

\VerdierSemanticAnchor{ch1-1-2-1-eq1}
\begingroup
\makeatletter
\tagsleft@true
\let\veqno\@@leqno
\makeatother
\begin{equation}
  \begin{cases}
    T(g)\colon\mathcal C_e\longrightarrow\mathcal C_e,\\[0.4em]
    T(e)=\text{the identity on }\mathcal C_e.
  \end{cases}
  \tag{1.2.1.1}
\end{equation}
\endgroup

\VerdierSemanticAnchor{ch1-1-2-1-transition}
For every pair $(g,h)$ of elements of $G$, there is [18] an isomorphism of
functors:

\VerdierSemanticAnchor{ch1-1-2-1-eq2}
\begingroup
\makeatletter
\tagsleft@true
\let\veqno\@@leqno
\makeatother
\begin{equation}
  \begin{cases}
    c(h,g)\colon T(hg)\longrightarrow T(g)T(h),\\[0.4em]
    c(g,e)=c(e,g)=\text{the identity on }T(g)
  \end{cases}
  \tag{1.2.1.2}
\end{equation}
\endgroup
